% ============================================================
% 192_T0_Algebraische_Kompositionsgrenzen_En_ch.tex
% Chapter content: Algebraic Composition Limits in FFGFT
% Why structural terms must not be inserted into foreign formulas
% Johann Pascher, May 2026
% ============================================================

\section*{Abstract}

In FFGFT, particle masses are derived from an algebraic torus structure:
$m_\ell = r_\ell \cdot \xi^{p_\ell} \cdot v$. These terms implicitly carry
the geometry of their origin. When inserted into formulas from a different
algebraic context — such as the Koide relation or Standard Model scattering
amplitudes — they generate structure-alien cross-terms that belong to neither
the FFGFT geometry nor the target formula. This error is not a numerical
problem but a category error: it violates the composition limits between
different algebraic structure spaces. This document analyses the error
systematically, shows its manifestations in FFGFT, and states the correct
rule for handling theory-internal terms in external formulas.

% ============================================================
\section{The Problem: Terms versus Results}
% ============================================================

\subsection{Two levels in FFGFT}

FFGFT operates on two clearly separated levels:

\begin{description}
  \item[Structure level (terms):] Symbolic expressions such as
    $m_\tau = r_\tau \cdot \xi^{2/3} \cdot v$ describe the
    \emph{geometric origin} of the mass. They encode the torus topology
    (exponent $p = 2/3$), the recursion depth (prefactor $r_\tau = 25/9$)
    and the scale (vacuum expectation value $v$).
  \item[Result level (numbers):] The numerical value
    $m_\tau \approx 1783\,\text{MeV}$ is a number without structural
    memory. It can be inserted like a measured value into any formula.
\end{description}

\textbf{The rule:} Results (numbers) may be inserted anywhere.
Terms (algebraic expressions) may only be used within the algebraic space
from which they originate.

\subsection{The Koide example}

The Koide relation reads:
\begin{equation}
  Q = \frac{m_e + m_\mu + m_\tau}{(\sqrt{m_e}+\sqrt{m_\mu}+\sqrt{m_\tau})^2}
  = \frac{2}{3}
  \label{eq:koide}
\end{equation}
This is experimentally confirmed to $0{.}001\,\%$ (PDG values).

\subsubsection*{Permitted use}

One evaluates the FFGFT terms numerically:
\begin{align*}
  m_e &\approx 0{.}511\,\text{MeV}, \\
  m_\mu &\approx 105{.}7\,\text{MeV}, \\
  m_\tau &\approx 1783\,\text{MeV},
\end{align*}
and inserts these numbers into \eqref{eq:koide}. One obtains
$Q_{\text{FFGFT}} \approx 0{.}6677 \neq 2/3$, a deviation of $0{.}16\,\%$.
This is correct and honest: the FFGFT masses are approximations
($\sim 1\,\%$ accuracy), and this imprecision propagates into $Q$.

\subsubsection*{Not permitted}

One inserts the FFGFT terms \emph{symbolically} into \eqref{eq:koide}:
\begin{equation}
  Q \stackrel{?}{=}
  \frac{r_e\,\xi^{3/2} + r_\mu\,\xi + r_\tau\,\xi^{2/3}}
  {(\sqrt{r_e\,\xi^{3/2}}+\sqrt{r_\mu\,\xi}+\sqrt{r_\tau\,\xi^{2/3}})^2}
  \stackrel{?}{=} \frac{2}{3}
  \label{eq:wrong}
\end{equation}
and attempts to verify this equation algebraically.

This fails. The reason is analysed in the next section.

% ============================================================
\section{Cross-term Analysis}
% ============================================================

\subsection{Where do the cross-terms come from?}

The denominator of $Q$ contains $(\sqrt{m_e}+\sqrt{m_\mu}+\sqrt{m_\tau})^2$.
Inserting the FFGFT terms:
\begin{align*}
  \sqrt{m_e} &\propto \xi^{3/4}, \\
  \sqrt{m_\mu} &\propto \xi^{1/2}, \\
  \sqrt{m_\tau} &\propto \xi^{1/3},
\end{align*}
squaring the denominator yields three \emph{diagonal terms}:
\begin{equation}
  m_e + m_\mu + m_\tau \propto \xi^{3/2} + \xi + \xi^{2/3}
\end{equation}
(the same powers as in the numerator) — and three \emph{cross-terms}:
\begin{align}
  2\sqrt{m_e\,m_\mu} &\propto \xi^{3/4+1/2} = \xi^{5/4}, \label{eq:k1}\\
  2\sqrt{m_e\,m_\tau} &\propto \xi^{3/4+1/3} = \xi^{13/12}, \label{eq:k2}\\
  2\sqrt{m_\mu\,m_\tau} &\propto \xi^{1/2+1/3} = \xi^{5/6}. \label{eq:k3}
\end{align}

\subsection{Why these cross-terms are structure-alien}

The exponents in the FFGFT structure space are multiples of $1/3$:
\[
  p \in \left\{\frac{2}{3},\; 1,\; \frac{3}{2},\; \frac{1}{3},\; \frac{1}{2},\; \frac{3}{4}\right\}.
\]
The cross-term exponents $5/4$, $13/12$, $5/6$ do \emph{not} belong to this
set. They arise only through the square-root operation of the Koide formula
and have no counterpart in the torus geometry.

\begin{keypoint}[Structure-alien terms]
Cross-terms from the Koide operation lie outside the FFGFT structure space.
They are neither geometrically interpretable nor physically meaningful.
They are an artefact of formula composition.
\end{keypoint}

\subsection{Condition for $Q = 2/3$ exactly}

For $Q = 2/3$ to hold algebraically exactly, one would need:
\begin{equation}
  m_e + m_\mu + m_\tau = 4\left(\sqrt{m_e\,m_\mu} + \sqrt{m_e\,m_\tau} + \sqrt{m_\mu\,m_\tau}\right).
  \label{eq:condition}
\end{equation}
This equation links the diagonal terms to the cross-terms.
It is not satisfied for $\xi = 4/30000$. The $\xi$-value for which it holds
exactly is $\xi_{\text{Koide}} \approx 1{.}372 \times 10^{-4}$ —
$2{.}9\,\%$ larger than the geometric value $\xi_{\text{par}} = 4/30000$.

This shows: the Koide relation enforces an algebraic constraint
that is independent of the FFGFT torus geometry.

% ============================================================
\section{The General Principle: Composition Limits}
% ============================================================

\subsection{Algebraic structure spaces}

Every physical theory defines an \emph{algebraic structure space}:
the set of expressions that are well-formed within the theory, and the
operations that are permitted within this space.

\begin{description}
  \item[FFGFT structure space:] Expressions of the form $r \cdot \xi^{p} \cdot v$
    with $p \in \frac{1}{3}\mathbb{Z}$. Permitted operations: multiplication,
    division, exponentiation by multiples of $1/3$.
  \item[Koide space:] Formula \eqref{eq:koide} with the operation
    $\sqrt{\cdot}$. This operation leaves the FFGFT structure space:
    $\sqrt{\xi^{2/3}} = \xi^{1/3}$ is still within it,
    but $\xi^{1/3} \cdot \xi^{1/2} = \xi^{5/6}$ is no longer in $\frac{1}{3}\mathbb{Z}$.
\end{description}

\begin{foundation}[Composition limit]
Two algebraic structure spaces $\mathcal{A}$ and $\mathcal{B}$ are
\emph{compositionally compatible} if all operations in $\mathcal{B}$
leave the space $\mathcal{A}$ invariant. Otherwise a composition limit
exists: terms from $\mathcal{A}$ must not be symbolically inserted
into formulas from $\mathcal{B}$.
\end{foundation}

\subsection{Further manifestations in FFGFT}

The Koide example is not unique. The same problem arises with:

\subsubsection*{Quantum numbers in Standard Model formulas}

FFGFT quantum numbers (torus winding numbers, $Z_3$ charges) have a
specific algebraic origin. They must not be inserted directly into
renormalisation equations or SM scattering amplitudes, as these formulas
presuppose different algebraic conventions.

\subsubsection*{Unit systems as a special case}

The most familiar example of composition limits is unit systems.
An expression in natural units ($\hbar = c = 1$) must not be inserted
without conversion into an SI formula. FFGFT terms carry an analogous
implicit ``unit'': the torus geometry.

\subsubsection*{Scale hierarchies}

The FFGFT terms $\xi^{p_\ell} \cdot v$ are calibrated for the particle sector.
They must not be inserted without scale correction into cosmological formulas,
where different levels of the $\xi$-hierarchy apply.

\subsection{The correct procedure}

\begin{keyresult}[Composition rule]
\textbf{Step 1:} Evaluate the FFGFT term numerically (result as a number).\\
\textbf{Step 2:} Insert that number like a measured value into the target formula.\\
\textbf{Step 3:} Honestly declare the result as an approximation
  when the term itself is only an approximation.
\end{keyresult}

This is structurally identical to error propagation in experimental physics:
before inserting a value into a new formula, one separates number and uncertainty.
The uncertainty of the FFGFT term ($\sim 1\,\%$) propagates — and produces
a different deviation in nonlinear formulas such as Koide than in linear ones.

% ============================================================
\section{Consequences for FFGFT Documentation}
% ============================================================

\subsection{What the Koide deviation means}

The deviation $Q_{\text{FFGFT}} \approx 0{.}6677 \neq 2/3$ is not an error
of the theory and not a contradiction of the experimental confirmation of
the Koide relation. It is the correct propagation of the FFGFT mass accuracy
($\sim 1\,\%$) through a nonlinear formula with a cross-term problem.

The experimental Koide relation $Q = 2/3$ confirms the FFGFT exponent
structure $\Delta p = 1/3$ by a different route: it shows that the
\emph{true} masses satisfy condition \eqref{eq:condition} — and that the
FFGFT terms approach this condition to $0{.}16\,\%$
without enforcing it.

\subsection{Linguistic consequence}

In external communications, a distinction must be made:

\begin{center}
\resizebox{\textwidth}{!}{%
\begin{tabular}{p{5.5cm}p{4cm}p{4cm}}
\toprule
\textbf{Statement} & \textbf{Status} & \textbf{Reason} \\
\midrule
``FFGFT explains the Koide formula through the exponent step $1/3$'' &
✅ Correct &
Structural statement about origin \\[6pt]
``FFGFT calculates $Q = 2/3$ exactly'' &
❌ Incorrect &
Composition limit violation \\[6pt]
``FFGFT masses agree with the Koide relation to $0{.}16\,\%$'' &
✅ Correct &
Honest numerical statement \\[6pt]
``The Koide relation is only valid with PDG values'' &
✅ Correct &
Context of the original formula \\
\bottomrule
\end{tabular}
}
\end{center}

% ============================================================

\section*{Connection to Doc.~193}

\textbf{Document~070} (\textit{On the Mathematical Structure of T0 Theory:
Why Products of Different $r_i$ Are Structurally Meaningless})
addresses the same fundamental principle at the level of mass formulas.

Doc.~193 proves the \textbf{non-closure theorem}: the set of admissible
pairs $(r_i, p_i)$ in FFGFT is not closed under multiplication:
\[
  \forall\, i \neq j:\quad (r_i \cdot r_j,\; p_i + p_j) \notin \mathcal{P}
\]
This means: $r_e \cdot r_\mu = 4/3 \cdot 16/5 = 64/15$ belongs to
no torus mode and no particle in FFGFT. The $r_i$ are not numbers in
a field but labels of a discrete geometric space --- their composition
is not ordinary field multiplication.

Doc.~193 shows: \emph{Mass formulas must not be inserted symbolically
into other formulas, because this generates cross-products $r_i \cdot r_j$
that lie outside $\mathcal{P}$.}

Doc.~192 (this document) generalises: \emph{Terms from one algebraic
structure space must not be symbolically inserted into formulas of a
foreign structure space} --- this applies to the Koide relation, to
quantum numbers in SM formulas, and to all scale hierarchies in FFGFT.

Both documents describe the same composition limit:
Doc.~193 at the level of mass formula terms (non-closure theorem),
Doc.~192 at the level of formula composition between different theories.

\section*{Summary}
% ============================================================

The composition limit between the FFGFT structure space and the Koide formula
explains the observed $0{.}16\,\%$ deviation in $Q$. It is not a theory error
but an indication of a categorical distinction: structural terms (with implicit
geometry) and numerical results (numbers without structural memory) must not
be used synonymously in foreign algebraic contexts.

The principle is general: it applies to quantum numbers in SM formulas,
to unit system transitions, and to scale hierarchies in FFGFT. The correct
procedure is always: evaluate first, then insert.
